\section{N-body}

Within $t_{2b}$, usually much shorter than the age of the universe for common galaxies, 
granularity is unimportant, collisions are negligible and stars can be described as
 particles moving in a smooth potential
generated by a continous $\rho$. Hence, stars will have a DF satisfying the CBE: 

\begin{equation}\label{CBE}
    \frac{Df}{Dt}=\frac{\partial f}{\partial t}+
    \mathbf{\dot{w}}\cdot \frac{\partial f}{\partial \mathbf{w}}=
    \frac{\partial f}{\partial t}+[f,H]=0
\end{equation}



In cartesian coordinates the Hamiltonian is
\begin{equation}\label{Hamiltonian in cartesian coord.}
    H(\mathbf{x},\mathbf{v})=\frac{1}{2}v^2+\Phi(\mathbf{x})
\end{equation}

and the CBE becomes

\begin{equation}\label{CBE in cartesian coord.}
    0=\frac{\partial f}{\partial t}+\mathbf{v}\cdot 
    \frac{\partial f}{\partial \mathbf{x}}-
    \frac{\partial \Phi}{\partial \mathbf{x}}
    \cdot \frac{\partial f}{\partial \mathbf{v}}
\end{equation}

Several times beyond $t_{2b}$ the system enters the collisional regime and first order 
corrections to the CBE are needed; the Fokker-Planck equation is an example of this:
\begin{equation}
    \frac{Df}{Dt}=C[f]
\end{equation}
where simplifications on the collision operator $C[f]$ can be done.\\In this scenario $f$ is 
no longer conserved along the orbit of the potential $\Phi$ and diffusion in the PS occurs, 
both $E$ and $L$ of a single star can change due to collisions, at odds with the CBE regime 
where, for istance $DH/Dt=0$. It is used in the loss cone problem.\\In order to exploit the 
properties of the CBE it is usefull to introduce the velocity moments, which will appear 
inside the Jeans equations:
\begin{equation}
    \overline{v_i}(\mathbf{x},t)\equiv 
    \frac{1}{\rho(\mathbf{x},t)}\int_\mathcal{R}^3 d^3 \mathbf{v} f v_i
\end{equation}
\begin{equation}
    \sigma_{ij}^2(\mathbf{x},t)\equiv 
    \frac{1}{\rho(\mathbf{x},t)}\int_\mathcal{R}^3 d^3 \mathbf{v}
    f (v_i-\overline{v_i})(v_j-\overline{v_j})=\overline{v_iv_j}-
    \overline{v_i}\,\overline{v_j}
\end{equation}
A function $I(\mathbf{x},\mathbf{v})$
is an integral of motion if and only if
\begin{equation}\label{Integral of motion def}
    \frac{D}{Dt}[I(\mathbf{x}(t),\mathbf{v}(t))]=0
\end{equation}
along any orbit of the potential $\Phi$.
\\Jeans Theorem: 
Any steady-state solution of the collisionless Boltzmann equation depends on 
the phase-space coordinates only through integrals of motion in the given potential, and any 
function of the integrals yields a steady-state solution of the collisionless Boltzmann 
equation.\\In a steady-state potential the Hamiltonian is an integral of motion; therefore $f$
will be any non-negative function of $H$.
\\Ergodic DFs: they are of the kind $f=f(H)$.
\\Equilibrium stellar systems have $\overline{v_i}=0$ and isotropic velocity-dispersion tensor
(being symmetric it can be diagonalized at any point in a proper orthogonal set of axes), i.e 
$\sigma^2_{ij}=\sigma^2 \delta_{ij}$. Thus every system with an ergodic DF is isotropic.
\\If the potential $\Phi$ is time-indipendent and spherically symmetric, then $\mathbf{L}$ 
appears as three isolating integral of motions but DF will depend on them only through 
$|\mathbf{L}|=L$, i.e. $f=f(H,L)$.Adopting spherical coordinates and splitting 
$v^2=v_r^2+v_t^2$, it holds 
\begin{equation}
    \overline{v_r}=0
\end{equation}
\begin{equation}
    \overline{\mathbf{v_t}}=\mathbf{0}
\end{equation}
and in general $\sigma_{ij}$ is diagonal with $\sigma^2_r\ne \sigma^2_\theta=\sigma^2_\phi$.\\Recovering
the DF from $\rho$ can be done by Eddington's formula, which applies to spherically symmetric,
globally isotropic stellar systems,
\begin{equation}\label{Eddington's formula}
f(\varepsilon)=\frac{1}{\sqrt{8}\pi^2}\frac{d}{d\varepsilon}
\int_{\varepsilon_t}^{\varepsilon}
\frac{d\rho}{d\Psi}\frac{d\Psi}{\sqrt{\varepsilon-\Psi}}
\end{equation}
with $\varepsilon=-H$ and $\Psi=-\Phi$ being the relative energy and potential respectively
and $\varepsilon_t$ the truncation energy of the system.\\Given the density profile, IC for N-body
simulations are generated by OpOpgadget software by sampling uniformly the internal mass profile
and the potential produced by the smooth profile $\rho$.\\The Sersic profile introduced by 
\textcite{Prugniel} has (with standard substitution $Z=b\left(\frac{r}{R_e}\right)^{1/n}$)
\begin{equation}\label{Sersic mass profile}
    M(<r)=4\pi \rho_0 R_e^3 n b^{n(p-3)}\gamma\left[n(3-p),b\left(\frac{r}{R_e}\right)^
    {1/n}\right]=M_*\frac{\gamma\left[n(3-p),b\left(\frac{r}{R_e}\right)^
    {1/n}\right]}{\Gamma[n(3-p)]}
\end{equation}
with 
\begin{equation}\label{complete gamma fun}
    \Gamma(a)=\int_{0}^{\infty}dt e^{-t}t^{a-1}
\end{equation}
the "Complete gamma function", 
\begin{equation}\label{upper incomplete gamma fun}
    \Gamma(a,x)=\int_{x}^{\infty}dt\,e^{-t}t^{a-1}
\end{equation}
the Upper incomplete gamma function and
\begin{equation}\label{lower incomplete gamma fun}
    \gamma(a,x)=\int_{0}^{x}dt\,e^{-t}t^{a-1}
\end{equation}
the Lower incomplete gamma function. By definition follows:
\begin{equation}\label{gamma fun relation}
    \Gamma(a)=\Gamma(a,x)+\gamma(a,x)
\end{equation}
The spatial distribution is spherical and untruncated but the mass profile converges to 
the finite value $M_*$, so the potential can be found from
\begin{equation}\label{spherical potential, general formula}
    \Phi(r)=\Phi(r_0)+\frac{GM(<r_0)}{r_0}-\frac{GM(<r)}{r}-4\pi G \int_{r}^{r_0}dt\,t \rho(t)
\end{equation}
Setting $r_0\rightarrow \infty$ and $\Phi(r_0)=0$, it becomes
\begin{equation}\label{spherical potential, simplified formula}
    \Phi(r)=-\frac{GM(<r)}{r}-4\pi G \int_{r}^{\infty}dt\,t \rho(t)
\end{equation}
and computation returns
\begin{equation}
\Phi(r)=
-4\pi G \rho_0 R_e^2 n\, b^{\,n(p-2)}
\left\{
\frac{b^{-n}}{r/R_e}\,
\gamma\!\left[n(3-p),Z\right]
+
\Gamma\!\left[n(2-p),Z\right]
\right\}
\end{equation}
or equivalently
\begin{equation}
    \Phi(r)=-\frac{GM_*}{r}\left\{\frac{\gamma\!\left[n(3-p),Z\right]}{\Gamma[n(3-p)]
    }+Z^n\frac{\Gamma[n(2-p),Z]}{\Gamma[n(3-p)]}\right\}
\end{equation}
The generation of IC for the dwarf galaxy NGC 2O5, hereafter taken as a $z=0$ nucleated dwarf
galaxy, is carried by OpOpgadget; the Sersic profile defined in this software is based on the 
one analyzed by \textcite{Neto_1999}, i.e.
\begin{equation}\label{Sersic profile Neto}
    \rho(r)=\rho_0\left(\frac{r}{a}\right)^{-p}e^{-(r/a)^\nu},
\end{equation}
where $n=\nu^{-1}$, $a$ is the Sersic profile scalelength and no $b(n)$ is present. The same
model can be described through several combinations of indipendent parameters; here, in order 
to use the effective radii values obtained from sky survy, we need to use the analytical
approximation 
\begin{equation}
    \ln{R_{eff}/a}=\frac{0.6950-\ln{\nu}}{\nu}-0.1789
\end{equation}

\begin{figure}
    \centering
    \includegraphics[width=0.5\linewidth]{NGC205_bulge_scatter.png}
    \caption{aaa}
    \label{NGC205_bulge_scatter}
\end{figure}

\begin{figure}
    \centering
    \includegraphics[width=0.5\linewidth]{NGC205_NSC_scatter.png}
    \caption{aaa}
    \label{NGC205_NSC_scatter}
\end{figure}

\begin{figure}
    \centering
    \includegraphics[width=0.5\linewidth]{NGC205_scatter.png}
    \caption{aaa}
    \label{NGC205_scatter}
\end{figure}


\begin{figure}[h]
\centering
 \begin{subfigure}{0.3\textwidth}
        \centering
        \includegraphics[width=\linewidth]{NGC205_bulge_space_hist.png}
        \caption{A}
    \end{subfigure}
    \hfill
    \begin{subfigure}{0.3\textwidth}
        \centering
        \includegraphics[width=\linewidth]{NGC205_NSC_space_hist.png}
        \caption{B}
    \end{subfigure}
    \hfill
    \begin{subfigure}{0.3\textwidth}
        \centering
        \includegraphics[width=\linewidth]{NGC205_space_hist.png}
        \caption{C}
    \end{subfigure}
    
    \caption{aaa}
    \label{3 NGC205_space_hist}
\end{figure}

\begin{figure}[h]
\centering
 \begin{subfigure}{0.3\textwidth}
        \centering
        \includegraphics[width=\linewidth]{NGC205_bulge_velocity_hist.png}
        \caption{A}
    \end{subfigure}
    \hfill
    \begin{subfigure}{0.3\textwidth}
        \centering
        \includegraphics[width=\linewidth]{NGC205_NSC_velocity_hist.png}
        \caption{B}
    \end{subfigure}
    \hfill
    \begin{subfigure}{0.3\textwidth}
        \centering
        \includegraphics[width=\linewidth]{NGC205_velocity_hist.png}
        \caption{C}
    \end{subfigure}
    
    \caption{aaa}
    \label{3 NGC205_velocity_hist}
\end{figure}

\begin{figure}[h]
\centering
 \begin{subfigure}{0.3\textwidth}
        \centering
        \includegraphics[width=\linewidth]{NGC205_bulge_density_Nbody.png}
        \caption{A}
    \end{subfigure}
    \hfill
    \begin{subfigure}{0.3\textwidth}
        \centering
        \includegraphics[width=\linewidth]{NGC205_NSC_density_Nbody.png}
        \caption{B}
    \end{subfigure}
    \hfill
    \begin{subfigure}{0.3\textwidth}
        \centering
        \includegraphics[width=\linewidth]{NGC205_NSC_density_Nbody.png}
        \caption{C}
    \end{subfigure}
    
    \caption{aaa}
    \label{3 NGC205_density_Nbody}
\end{figure}

\begin{figure}[h]
\centering

\begin{subfigure}{0.45\textwidth}
\centering
\includegraphics[width=\linewidth]{NGC205_density_Nbody_vs_analytic.png}
\caption{A}
\end{subfigure}
\hfill
\begin{subfigure}{0.45\textwidth}
\centering
\includegraphics[width=\linewidth]{NGC205_mass_Nbody_vs_analytic.png}
\caption{B}
\end{subfigure}

\caption{aaa}
\label{NGC205 density and mass plot Nbody vs analytical}
\end{figure}

COM data:
\begin{itemize}
    \item position:
    \item velocity:
\end{itemize}

To test the stability of the Nbody generated IC, we let the system free to evolve over several
tens of the dynamical time $t_{dyn}$ at half-mass radius $r_{half}$; by definition, the former
is
\begin{equation}\label{dynamical time def}
    t_{dyn}=\frac{1}{\sqrt{G\bar{\rho}}}
\end{equation}

and the latter reads, for a spherical system,
\begin{equation}
    \bar{\rho}=\frac{M/2}{\frac{4}{3}\pi r_{half}^3}
\end{equation}

Numerically is found $r_{half}=...=... R_{eff}$.

\begin{figure}
    \centering
    \includegraphics[width=0.5\linewidth]{NGC205_halfmass.png}
    \caption{aaa}
    \label{NGC205_halfmass}
\end{figure}
Yet, we get $t_{dyn}=...$, which agrees within $x\%$ the estimate obtained from the Analysis
tool provided by OpOpgadget, which measures this quantity directely from the generated model.
\\Waiting for evolution with Ketju.\\Additionaly, the virial ratio $Q$ has been computed as well as it gives a first indication
of whether the model is virialized, i.e. at equilibrium, or expansion/contraction operate 
within several dynamical times. In $S_{CM}$, it is defined as
\begin{equation}\label{virial ratio}
    Q=\frac{2K}{|U|}
\end{equation}
where
\begin{equation}
K = \sum_{i=1}^{N} \frac{1}{2} m_i v_i^2
\end{equation}
and
\begin{equation}
U = -\frac{1}{2} \sum_{i=1}^{N} \sum_{j \ne i} \frac{Gm_i m_j}{|\mathbf{x_i}-\mathbf{x_j}|}
\end{equation}
stand for the total kinetic and gravitational potential energies respsectively.
\\By erasing the offset due the COM kinetic energy contribute, i.e.
\begin{equation}\label{virial ratio without offset}
    Q=\frac{2(K-K_{CM})}{|U|}
\end{equation}
it is found $Q=x$.